\section*{Plain Language Summary}

Archaeologists studying the plains of southern Iraq look for ancient mounds,
called tells, that mark the sites of settlements abandoned thousands of years
ago. Because these mounds are visible in satellite photographs, researchers
have started training computer programs to spot them automatically, saving
archaeologists from having to scan enormous stacks of imagery by eye. Training
such a program means feeding it thousands of small image files, over and over,
on a shared supercomputer built for exactly this kind of large computation.

The intuitive assumption is that the expensive graphics chip doing the actual
learning is the part that takes the most time, so that is where effort to speed
things up usually goes. This thesis shows that assumption is often wrong for
this kind of work. The real slowdown was frequently the plumbing: fetching
thousands of individual image files, one at a time, from a storage system that
is shared by everyone on the supercomputer at once. The expensive chip sat
idle, waiting for files, far more often than it sat busy computing.

To address this, the thesis builds a smarter file-loading system that watches
its own behaviour for a few seconds at the start of a training run and decides,
automatically, how best to supply data for the rest of that run: read files one
by one as before, bundle them into larger packages that are faster to read
sequentially, or copy the whole dataset onto fast local storage ahead of time.
Because the order in which files will be needed is known in advance, the system
can prepare data before it is asked for, rather than reacting after the fact.

Testing this system across three different AI models revealed two things. The
underlying idea works: when the smarter file-handling was actually used, the
expensive chip sat idle far less often, sometimes cutting waiting time by more
than half. But the automatic decision-maker itself was miscalibrated: it made
its judgement from a handful of very first measurements taken at the start of a
run, before some one-off start-up costs had settled down, and those inflated
measurements made ordinary file loading look faster than it really was. As a
result, the system kept choosing the plain, unoptimised approach almost every
time, even when a better one was available and demonstrably worked. This
thesis traces that miscalibration to its exact cause and shows that simply
ignoring the first few, skewed measurements restores accurate decision-making
at essentially no extra cost.

The broader lesson extends beyond archaeology: on shared supercomputers,
training large AI models is frequently limited by how efficiently data can be
delivered rather than by the model itself, and any system that tries to tune
itself automatically must be built carefully enough not to be fooled by its own
start-up noise.
